
\section{Marco referencial}

\subsection{Marco teórico}

Las herramientas digitales han adquirido relevancia en los procesos de enseñanza y aprendizaje de la educación superior, al facilitar el acceso a información, la comunicación entre los participantes, la elaboración de contenidos y el desarrollo de actividades académicas. Su incorporación, sin embargo, no garantiza por sí misma una mejora del aprendizaje, pues sus resultados dependen de las condiciones de acceso, las competencias de los usuarios, los objetivos pedagógicos y las formas de utilización.

El Informe de seguimiento de la educación en el mundo de la UNESCO examina las oportunidades y los desafíos de la tecnología en la educación y destaca la importancia de evaluar su pertinencia, equidad y contribución a los objetivos educativos \parencite{unesco2023tecnologia}. Desde esta perspectiva, estudiar el uso de herramientas digitales requiere considerar no solamente su disponibilidad, sino también la frecuencia con la que se utilizan, las actividades para las que se emplean y las dificultades que experimentan los estudiantes.

Para esta investigación, el análisis se orienta hacia una caracterización descriptiva de las prácticas digitales en el contexto universitario. No se parte del supuesto de que utilizar más herramientas implique necesariamente obtener mejores resultados académicos; en cambio, se busca identificar patrones de uso y percepciones que permitan comprender cómo se integran estas tecnologías en las actividades de estudio.

\subsection{Marco conceptual}

Para delimitar el objeto de estudio, se establecen los siguientes conceptos operativos:

\begin{itemize}
    \item \textbf{Herramientas digitales:} aplicaciones, plataformas y servicios tecnológicos que los estudiantes utilizan para desarrollar actividades académicas, como consultar información, elaborar documentos, comunicarse, colaborar o gestionar tareas.

    \item \textbf{Uso académico:} utilización de una herramienta digital con una finalidad vinculada con las actividades de aprendizaje, la elaboración de trabajos, la consulta de fuentes, la comunicación académica o la organización del estudio.

    \item \textbf{Frecuencia de uso:} periodicidad con la que un estudiante utiliza una herramienta para fines académicos. Se medirá mediante categorías ordenadas de respuesta, desde nunca hasta diariamente.

    \item \textbf{Finalidad de uso:} actividad concreta para la cual se emplea una herramienta digital. Se contemplan, entre otras, la búsqueda de información, la redacción, la comunicación, la colaboración y la organización de tareas.

    \item \textbf{Utilidad percibida:} valoración subjetiva que realiza el estudiante sobre la contribución de una herramienta a sus actividades académicas. Se recogerá mediante una escala de respuesta definida en el cuestionario.

    \item \textbf{Barreras de acceso y uso:} dificultades que pueden limitar la utilización de herramientas digitales, como la conectividad, la disponibilidad de dispositivos, los costos, las habilidades digitales, la accesibilidad o las preocupaciones relacionadas con la privacidad.
\end{itemize}

Estas definiciones tienen una finalidad operativa: orientar la construcción del cuestionario, establecer categorías de respuesta y mantener la correspondencia entre los objetivos, las variables y el análisis de los datos. Las categorías definitivas deberán revisarse antes de la aplicación del instrumento.

\subsection{Marco contextual}

La investigación se delimita geográficamente a Bogotá, Distrito Capital, Colombia, y se dirige a estudiantes de 18 años o más matriculados en programas de educación superior de instituciones ubicadas en la ciudad durante el periodo que se establezca para la recolección de información.

Como referencia contextual sobre las condiciones de conectividad, el Departamento Administrativo Nacional de Estadística reportó que, en 2024, el porcentaje de hogares con conexión a internet fue del 82,7\,\% en Bogotá y del 65,6\,\% en el total nacional \parencite{dane2025entic}. Estas cifras permiten contextualizar las condiciones generales de acceso a internet, pero corresponden a hogares y no a estudiantes universitarios. Por tanto, no deben interpretarse como una medición directa de la conectividad individual, de la frecuencia de uso de herramientas digitales ni de las competencias tecnológicas de la población estudiada.

La selección de Bogotá como ámbito territorial permite delimitar la población objetivo y orientar la planificación del trabajo de campo. No obstante, la población accesible dependerá de las instituciones que autoricen el estudio, de los programas académicos incluidos y de la disponibilidad de un marco muestral. Estas condiciones deberán especificarse antes de determinar el tamaño de la muestra y el alcance de las conclusiones.

\subsection{Marco legal}

El proyecto considera, en primer lugar, la Ley 30 de 1992, mediante la cual se organiza el servicio público de la educación superior en Colombia \parencite{colombia1992ley30}. Esta norma constituye un referente general para situar la investigación en el ámbito de la educación superior; su mención no implica que regule específicamente el uso de las herramientas digitales analizadas.

En materia de información personal, se considera la Ley 1581 de 2012, que establece disposiciones generales para la protección de datos personales en Colombia \parencite{colombia2012ley1581}. En consecuencia, el diseño de la encuesta deberá aplicar criterios de minimización de datos, información clara a los participantes, acceso restringido a los registros y presentación agregada de los resultados.

La aplicación concreta de estas disposiciones dependerá de las características del tratamiento de datos y de los procedimientos institucionales aplicables. Antes de iniciar la recolección, se deberán verificar los requisitos de autorización, consentimiento, custodia y conservación de la información, así como las orientaciones éticas de las instituciones participantes.

\subsection{Integración del marco referencial}

Los componentes del marco referencial cumplen funciones complementarias. El marco teórico aporta una perspectiva para comprender las oportunidades y limitaciones de la tecnología en la educación; el conceptual define las variables que serán observadas; el contextual delimita el territorio y presenta información general sobre conectividad; y el legal orienta el tratamiento responsable de la información recopilada.

En conjunto, estos elementos sustentan una investigación descriptiva de enfoque cuantitativo, con diseño no experimental y transversal. Su función es orientar la recolección y la interpretación de los datos, sin sustituir los resultados empíricos que se obtendrán durante el trabajo de campo.


\begin{figure}[H]
    \caption{Delimitación contextual del estudio}
    \label{fig:delimitacion-contextual}

    \centering

    \begin{tikzpicture}[
        bloque/.style={
            draw=black!60,
            rounded corners=2pt,
            align=center,
            text width=5cm,
            minimum height=1.15cm,
            inner sep=8pt,
            font=\small
        },
        flecha/.style={
            -{Stealth[length=2.5mm]},
            thick,
            black!65
        }
    ]

        % Delimitación geográfica y poblacional
        \node[bloque] (ciudad) {
            \textbf{Bogotá, D. C.}
        };

        \node[bloque, below=0.35cm of ciudad] (instituciones) {
            Instituciones de educación superior
        };

        \node[bloque, below=0.35cm of instituciones] (poblacion) {
            Población propuesta: estudiantes de 18 años o más
        };

        % Objeto de estudio
        \node[
            bloque,
            text width=4.2cm,
            minimum height=2.2cm,
            right=1.3cm of instituciones
        ] (objeto) {
            \textbf{Objeto de estudio}\\[4pt]
            Herramientas digitales para fines académicos
        };

        % Relaciones conceptuales
        \draw[flecha] (ciudad) -- (instituciones);
        \draw[flecha] (instituciones) -- (poblacion);
        \draw[flecha] (instituciones.east) -- (objeto.west);

    \end{tikzpicture}

    \par\smallskip
    \begin{minipage}{0.9\textwidth}
        \footnotesize
        \textit{Nota.} Esquema conceptual de delimitación, no mapa cartográfico ni representación de la ubicación exacta de instituciones.
    \end{minipage}
\end{figure}

\newpage